\subsection{Incidence cubique et réalisation lorentzienne}

La complétion cubique apporte une relation entre intérieur et frontière.
Sa géométrie peut être développée au moyen de l'incidence des six faces
avec les huit sommets. Étiquetons les faces par \((i,\varepsilon)\),
\(i\in\{1,2,3\}\), \(\varepsilon\in\{+1,-1\}\), et les sommets par
\(\nu\in\{+1,-1\}^3\). La matrice d'incidence est
\[
 M_{(i,\varepsilon),\nu}=\mathbf1_{\{\nu_i=\varepsilon\}}.
\]
Soient \(u=(1,\ldots,1)^{\mathsf T}\in\RR^6\), \(R_F\) l'opposition des faces et \(P_u=uu^{\mathsf T}/6\), \(P_-=(I-R_F)/2\).

\begin{proposition}[Quotient d'incidence de dimension quatre]
Le gramien d'incidence satisfait
\[
 MM^{\mathsf T}=12P_u+4P_-.
\]
Ses valeurs propres \(12,4,0\) ont les multiplicités \(1,3,2\).
En conséquence,
\[
 Q_\square=\RR^6/\ker(MM^{\mathsf T})
 \simeq\RR u\oplus E_-,\qquad E_-=\operatorname{im}P_-.
\]
\end{proposition}
\begin{proof}
Une face contient quatre sommets ; deux faces opposées ont une intersection vide, et deux faces distinctes non opposées partagent deux sommets. La matrice \(12P_u+4P_-\) possède précisément ces entrées. L'espace uniforme est de dimension un ; les différences entre faces opposées engendrent un espace de dimension trois, orthogonal à l'espace uniforme. Son complément est de dimension deux et appartient au noyau.
\end{proof}

La décomposition \(1+3\) procède de l'incidence. Le choix temporel
est déclaré au moyen de la forme bilinéaire
\[
 B_\square([x],[y])=\langle x,(P_--P_u)y\rangle.
\]
Elle est bien définie sur le quotient, car les deux projecteurs annulent le
noyau. Dans la base
\[
 t=[u]/\sqrt6,\qquad
 d_i=[e_{i,+}-e_{i,-}]/\sqrt2,
\]
sa matrice est \(\operatorname{diag}(-1,1,1,1)\). Sont ainsi explicités
le domaine, la décomposition et la convention de signe qui produisent
la réalisation de Minkowski. Le gramien positif détermine le quotient ;
l'affectation temporelle détermine la forme indéfinie sur celui-ci.

L'opérateur
\[
 W_\square=\sqrt6P_u+\sqrt2P_-,\qquad
 W_\square e_{i,\varepsilon}=t+\varepsilon d_i
\]
envoie les six étiquettes de face sur des directions nulles, puisque
\(B_\square(t+\varepsilon d_i,t+\varepsilon d_i)=-1+1=0\).
L'inversion d'une direction et l'opposition des faces disposent ainsi
d'une réalisation géométrique concrète.

\subsubsection{Soudure, inversion et subdivision}

La réalisation cellulaire utilise un transport inversible \(U_m(a)\),
qui conserve la forme lorentzienne, entre les fibres aux extrémités d'une
arête duale orientée \(a\), une étiquette
de face \(\lambda_m(a)\) et une orientation \(\sigma_m(a)\).
Avec \(\ell_m=\ell_0/9^m\), l'application de soudure est
\[
 \beta_m(a)=\ell_m\sigma_m(a)
 W_{\square,s(a)}e_{\lambda_m(a)}.
\]
L'inversion compatible satisfait
\[
 \beta_m(\bar a)=-U_m(a)\beta_m(a).
\]
Identifions les fibres de l'origine à deux niveaux consécutifs au moyen de l'isométrie de raffinement déclarée. Si les neuf sous-segments \(a_1,\ldots,a_9\), de niveau \(m+1\), conservent leur étiquette et leur orientation lorsqu'ils sont transportés vers cette fibre commune, alors
\begin{equation}
 \beta_m(a)=\sum_{j=1}^9
 U_{m+1}(a_1\cdots a_{j-1})^{-1}\beta_{m+1}(a_j).
 \label{eq:soldadura-refinamiento}
\end{equation}
Chaque terme transporté équivaut à \(\ell_m/9\) suivant la même
direction nulle, ce qui démontre l'identité. L'hypothèse de
compatibilité détermine quelles subdivisions réalisent ce raffinement.
L'équation réunit direction, échelle et transport, au lieu de limiter
la correspondance géométrique à un dénombrement des dimensions.

L'orientation et la forme lorentzienne étant également fixées, la dualité de
Hodge dans \(\Lambda^2Q_\square^*\) satisfait \(\star^2=-I\).
En dimension quatre, le signe est \((-1)^{2(4-2)+1}=-1\).
On obtient une structure complexe réelle sur les deux-formes et,
après complexification, les sous-espaces propres de valeurs \(+i\) et
\(-i\). L'application électronique de l'Article I construit sa section
persistante sur le couple de feuillets APP et fixe la polarité qui intervient dans
la réalisation électromagnétique. Dans cet article, l'orientation et la forme
lorentzienne déclarées transportent cette structure vers le domaine gravitationnel :
la dualité de Hodge organise les composantes du courant, et la section
\ref{vii:sec:corriente-respuesta-cartan} les incorpore à l'équation
algébrique de Cartan--Holst.

\subsection{Compatibilité entre les réductions modulaire et positionnelle}

La base décimale par blocs intervient au moyen d'une division euclidienne
qui conserve simultanément le résidu et le quotient :
\[
 n=1000q+r,\qquad 0\le r<1000.
\]
Puisque \(1000\equiv1\pmod9\) et \(1000\equiv1\pmod3\),
\[
 n\equiv q+r\pmod9,\qquad n\equiv q+r\pmod3.
\]
Par conséquent, la réduction d'un bloc requiert la contribution du
quotient pour restituer la classe de l'entier. Les nombres \(0\) et
\(1000\) ont le même reste modulo \(1000\), mais des classes différentes
modulo \(9\) et modulo \(3\). C'est une raison opératoire de conserver
la mémoire pendant le changement de résolution.

Le raffinement ternaire et la détermination d'un préfixe décimal
peuvent se coordonner au moyen d'un unique résidu entier. Si \(N/3^L\)
est une extrémité de cylindre et \(D/1000^r\) une extrémité décimale, nous définissons
\[
 A=1000^rN-D3^L.
\]
En ajoutant six symboles ternaires, \(N'=729N+\nu\),
\(L'=L+6\), avec \(0\le\nu<729\). Si le préfixe décimal reste
fixe, on obtient
\[
 A'=729A+\nu1000^r.
\]
Si l'on incorpore en outre un bloc décimal \(\delta\), alors \(D'=1000D+\delta\), \(r'=r+1\), et
\begin{equation}
 A'=729000A+\nu1000^{r+1}-729\delta3^L.
 \label{eq:residuo-dos-bases}
\end{equation}
Les deux identités se démontrent en substituant les définitions. La
comparaison des extrémités se réduit ainsi à des opérations entières qui
retiennent la contribution de chaque prolongation. Cette compatibilité
permet de raffiner une région avant d'évaluer sa coordonnée réelle.
